% together/hazptr.tex
% mainfile: ../perfbook.tex
% SPDX-License-Identifier: CC-BY-SA-3.0

\section{风险指针助手}
\label{sec:together:Hazard-Pointer Helpers}
%
\epigraph{重要的是小事情，成百上千的小事情。}
	 {Cliff Shaw}

本节探讨一些可以借助 hazard pointer 解决的问题。
此外，hazard pointer 有时也可以用来解决
\cref{sec:together:RCU Rescues} 中提出的问题，反之亦然。

\subsection{可扩展引用计数}
\label{sec:together:Scalable Reference Count}

假设引用计数正在成为性能或可扩展性的瓶颈。
你能做什么？

一种方法是改用\IXpl{hazard pointer}。

两者之间存在一些差异，
最显著的也许是使用 hazard pointer 时，
确定相应的引用计数何时达到零的代价极其昂贵。

解决此问题的一种方法是在引用计数和 hazard pointer 之间分担负载。
每个数据元素都有一个引用计数器，
一方面跟踪引用该元素的其他数据元素的数量，
另一方面读者使用 hazard pointer。

要使这种安排既高效又正确可能相当具有挑战性，
因此有兴趣的读者可以查看
Folly 开源库中实现的 UnboundedQueue 和 ConcurrentHashMap 数据结构。\footnote{
	\url{https://github.com/facebook/folly}}

\subsection{长时间访问}
\label{sec:together:Long-Duration Accesses}

假设一个读写锁的读者持有锁的时间太长，
导致更新被过度延迟。
如果该读者可以合理地转换为使用引用计数
而非读写锁，但性能和可扩展性的考虑又阻止使用
实际的引用计数器，那么 hazard pointer 提供了
引用计数的一种可扩展变体。

关键点在于，读写锁的读者会阻塞该锁的所有更新，
而 hazard pointer 只是保持住实际需要的数据，
同时仍允许更新继续进行。

如果该读者无法合理地转换为使用引用计数，
\cref{sec:together:Long-Duration Accesses Two}
中的技巧可能会有所帮助。

% @@@ papers to maybe cite: OrcGC, ThreadScan, Fast and Robust Memory...

% @@@ Generalized hazard-pointer link counts, if and when.

% @@@ Representative hazard pointer for list, so that nothing
% @@@ in list gets freed until list's hazard pointer is released.
% @@@ Midpoint between hazard pointers and RCU, in fact, you
% @@@ could argue that Tasks Trace RCU with read-side memory
% @@@ barriers is sort of a per-CPU hazard pointers implementing RCU.
% @@@ Except no re-checking because CPUs cannot be freed.
